% Section 3.3, from lectures/L03/appendix/b_exercises.md.

\section{Exercises}
\label{c3:sec:exercises}\label{c3:app:b}

\begin{aside}
\textbf{How to work these.} The Reflection and Calculation exercises are answered on paper, from
the chapter. The Code exercises are programs for the board: build and flash them the way
\secref{c2:app:c} (Microchip Studio) or \secref{c2:app:d} (Linux) describes, or run them in the
simulator of Appendix~\ref{sim:ch} if you have no board yet.

\textbf{How to check your work.} The Reflection and Calculation exercises are answered in
Appendix~\ref{app:answers}. The Code exercises have no printed solution: the board tells you
whether a program does what the exercise describes. Try each exercise before you read its answer.
\end{aside}

% ------------------------------------------------------------------------------------------
\exerciseset{Concepts}
\label{c3:set:1}

\reflexercise{Polling versus interrupts}
\label{c3:ex:1}
In your own words (2--3 sentences), explain why polling \code{BTN1_PRESSED} in a \code{while} loop
scales poorly once you have several inputs to watch, and how interrupts solve that.

\reflexercise{\code{EICRA} and \code{EIMSK}}
\label{c3:ex:2}
What does \code{EICRA} configure that \code{EIMSK} doesn't? Could you enable INT0 without touching
\code{EICRA} at all; what would happen if you did?

% ------------------------------------------------------------------------------------------
\exerciseset{INT0/INT1}
\label{c3:set:2}

\codeexercise{Moving the button to INT1}
\label{c3:ex:3}
Rewire the INT0 example of \secref{c3:a:3} (button on D2) to use \textbf{INT1} on D3 instead.
Write out the three configuration lines that change.

\reflexercise{The wrong edge}
\label{c3:ex:4}
A colleague configures \code{EICRA} for a falling edge instead of rising, but leaves the rest of
the INT0 example unchanged, and complains the LED toggles when the button is \emph{released}, not
pressed. Explain what happened.

% ------------------------------------------------------------------------------------------
\exerciseset{Pin Change Interrupts}
\label{c3:set:3}

\calcexercise{An interrupt on D4}
\label{c3:ex:5}
A project needs an interrupt on D4. Work out, from the formula table in \secref{c3:a:5}, which
\code{PCINT} number that is, which port group it belongs to, and which three registers/bits need to
be touched to enable it.

\reflexercise{Two buttons, one vector}
\label{c3:ex:6}
Two buttons are wired to D11 and D12 (both \code{PORTB}). Both are configured via PCI and share
\code{PCINT0_vect}. Sketch how a single ISR could tell which button changed, using a ``previous
state'' byte compared against the current \code{PINB} read.

% ------------------------------------------------------------------------------------------
\exerciseset{Program}
\label{c3:set:4}

\codeexercise{A PCI blink toggle}
\label{c3:ex:7}
Wire up an LED on D6 and a button on D4. Using a Pin Change Interrupt (not INT0/INT1) on the
button, write a complete program, from scratch, that:
\begin{itemize}
  \item Starts with the LED off and blinking disabled.
  \item Each button press toggles between two states: blinking (LED on/off every 200\,ms) and
    disabled (LED held off).
  \item The ISR only flags that the button was pressed, filtering the release edge the same way as
    the PCI example of \secref{c3:a:5}; the actual state toggle and the blink timing both happen in
    \code{main()}.
\end{itemize}
You'll need a blocking delay for the blink timing itself (\code{_delay_ms()}, as in
Chapter~\ref{c2:ch}); non-blocking delays are the subject of Chapter~\ref{c5:ch}, on timers, not
of this one. This chapter's own \secref{c3:a:4} (blink an LED via INT1) is the closest worked
example to lean on.

\codeexercise{An interrupt-driven state machine}
\label{c3:ex:8}
Build a 5-state state machine, driven entirely by interrupts, using every kind covered in this
chapter at once.

Wire up three LEDs, \code{LED1} on D8, \code{LED2} on D9, and \code{LED3} on D10, plus five
buttons. Declare the five states as an \code{enum} (\code{STATE_OFF}, \code{STATE_SLOW},
\code{STATE_MEDIUM}, \code{STATE_FAST}, \code{STATE_ON}), and wire each button as follows,
advancing/rewinding with wraparound at both ends of the range:
\begin{itemize}
  \item \code{BTN1} on D2 (\textbf{INT0}, falling edge) is a reset: jump straight to
    \code{STATE_OFF}, ignoring the normal advance/rewind stepping.
  \item \code{BTN2} on D3 (\textbf{INT1}, rising edge) rewinds the state by two steps.
  \item \code{BTN3} on D11 (\textbf{PCI}) rewinds the state by one step.
  \item \code{BTN4} on D12 (\textbf{PCI}) advances the state by one step.
  \item \code{BTN5} on D13 (\textbf{PCI}) advances the state by two steps.
\end{itemize}
The three LEDs' behavior depends on the current state:
\begin{itemize}
  \item \code{STATE_OFF}: all three LEDs held off.
  \item \code{STATE_SLOW} / \code{STATE_MEDIUM} / \code{STATE_FAST}: the three LEDs chase in a loop
    (\code{LED1} $\rightarrow$ \code{LED2} $\rightarrow$ \code{LED3} $\rightarrow$ \code{LED1}
    \ldots), lighting exactly one at a time, holding each for 500\,ms / 250\,ms / 50\,ms
    respectively.
  \item \code{STATE_ON}: all three LEDs held on.
\end{itemize}
As with every other example in this chapter, keep each ISR minimal: flag which button fired (and,
for the three PCI buttons sharing one vector, work out \emph{which} pin changed, as in
Exercise~\ref{c3:ex:6}), and let \code{main()} do the actual work of updating \code{state} and
driving the LEDs.

\bookimage{lectures/L03/appendix/images/circuit_state_machine_exercise.jpg}{Three LEDs on D8--D10
and five buttons, on D2, D3 and D11--D13.}{c3:fig:statemachine}
